\documentclass[12pt,a4paper]{article}

% ===== بسته‌های اصلی =====
\usepackage{xepersian}
\settextfont{Estedad-Medium.ttf}              
\setlatintextfont{Latin Modern Roman} 


\usepackage{amsmath}
\usepackage{amssymb}
\usepackage{geometry}
\usepackage{booktabs}
\usepackage{array}
\usepackage[hidelinks]{hyperref}

\geometry{margin=2.5cm}
\linespread{1.3}

\title{گزارش پروژه‌ی داده‌کاوی}
\author{}
\date{}

\begin{document}

\begin{center}
    {\Large \textbf{گزارش پروژه‌ی داده‌کاوی}}\\[6pt]
    \begin{latin}
        {\large \textbf{Semantic Duplicate \& Near-Plagiarism Detection Engine}}
    \end{latin}\\[10pt]
    \textbf{دانشجویان:} مریم صادقی، شیرین شفیعی \\
    \textbf{استاد درس:} دکتر شاکری
\end{center}

\noindent\hrulefill

\section{توضیح روش}

در این پروژه یک سامانه‌ی خط‌فرمان (CLI) برای تشخیص اسناد مشابه، کپی‌شده یا بازنویسی‌شده پیاده‌سازی شده است. دو مسیر اصلی مورد مقایسه قرار گرفتند:

\begin{itemize}
    \item \textbf{مسیر اول:} Shingling + MinHash + LSH (تخمین شباهت جاکارد با استفاده از امضای فشرده و باندبندی)
    \item \textbf{مسیر دوم:} TF-IDF + SimHash (وزن‌دهی معنایی و فاصله‌ی همینگ)
\end{itemize}

در هر دو مسیر از یک روش پیش‌پردازشی استفاده شده است. ما پیش‌پردازش متون (پاک‌سازی، توکن‌سازی، حذف ایست‌واژه و شینگل‌سازی) را در دو قسمت انجام دادیم، ابتدا همه مراحل را انجام دادیم به غیر از شینگل‌سازی چون هر دو الگوریتم ما به توکن‌ها نیاز داشتند ولی الگوریتم مورد استفاده در مسیر دوم به شینگل‌سازی نیازی نداشت بنابراین ما اولین مرحله پیش‌پردازش را انجام دادیم که داده برای استفاده هر دو الگوریتم مناسب باشد، سپس خروجی این مرحله را در فایل‌های JSON ذخیره کردیم (توکن‌ها را مرتب شده در این فایل ذخیره کردیم تا برای شینگل‌سازی مشکلی پیش نیاید) و آخرین مرحله پیش‌پردازش را برای الگوریتم اولی پیاده کردیم.

\subsection{پیش‌پردازش (Preprocessing)}

پیش‌پردازش مطابق با دستورالعمل پروژه و فایل «متن‌کاوی» انجام شده است:

\begin{enumerate}
    \item \textbf{پاک‌سازی:}
    \begin{itemize}
        \item lowercasing
        \item حذف علائم نگارشی و کاراکترهای غیرمجاز با نگه‌داری \texttt{-} و \texttt{'} (برای کلمات مرکب و مخفف)
        \item حذف فاصله‌های اضافی
        \item تشخیص متون غیرانگلیسی (با آستانه‌ی ۷۰٪ کاراکترهای مجاز) و رد آن‌ها
    \end{itemize}
    \item \textbf{توکن‌سازی:} با استفاده از \texttt{nltk.word\_tokenize} و فیلتر کردن توکن‌های purely-punctuation
    \item \textbf{حذف ایست‌واژه:} با استفاده از لیست استاندارد NLTK (انگلیسی) و فیلتر کردن توکن‌های purely-punctuation
    \item \textbf{Stemming:} با الگوریتم Porter (برای کاهش کلمات به ریشه)
    \item \textbf{شینگل‌سازی:}
    \begin{itemize}
        \item در حالت عادی: شینگل‌های کلمه‌ای با اندازه‌ی \texttt{k=3} (word n-grams)
        \item در حالت fallback (متون کوتاه): شینگل‌های کاراکتری با همان \texttt{k}
    \end{itemize}
\end{enumerate}

خروجی مرحله‌ی پیش‌پردازش، \textbf{لیست توکن‌ها} (به‌همان ترتیب) است که در پوشه‌ی \texttt{data/processed/} به‌صورت فایل‌های JSON ذخیره می‌شود.

\subsection{مسیر اول: MinHash + LSH}

\subsubsection*{الف) MinHash}

\begin{itemize}
    \item \textbf{تعداد توابع هش:} \texttt{num\_perm = 128}
    \item \textbf{تابع هش پایه:} $h(x) = (a \cdot x + b) \bmod p$ با $p = 2^{61} - 1$
    \item \textbf{تبدیل شینگل به عدد:} \texttt{zlib.crc32} (سریع و با توزیع یکنواخت)
    \item \textbf{امضا:} بردار ۱۲۸-تایی از مقادیر مینیمم هر تابع هش
    \item \textbf{تخمین شباهت جاکارد:} نسبت برابری مؤلفه‌های متناظر دو امضا
\end{itemize}

\subsubsection*{ب) LSH (باندبندی)}

\begin{itemize}
    \item \textbf{پارامترها:} \texttt{bands = 64}, \texttt{rows = 2} (زیرا $64 \times 2 = 128$)
    \item \textbf{تابع هش باند:} $H(v_1,\dots,v_r) = (a_1 v_1 + \cdots + a_r v_r + b) \bmod p$ با $p = 2^{31} - 1$
    \item \textbf{تولید جفت‌های کاندید:} دو سندی که در حداقل یک باند در یک سطل مشترک باشند، کاندید محسوب می‌شوند.
    \item \textbf{محاسبه‌ی شباهت واقعی:} برای هر جفت کاندید، جاکارد دقیق (از روی شینگل‌ها) محاسبه شده و در صورت عبور از آستانه (که با استفاده از منحنی S به‌دست آمده) به‌عنوان تکراری تأیید می‌شود.
\end{itemize}

یک نکته‌ای است که ما ابتدا مقادیر را این‌گونه در نظر گرفته بودیم، \texttt{rows = 4} و \texttt{bands = 32}، ولی با توجه به اینکه از دیتاست PAN برای پروژه استفاده کردیم توی خیلی از اسناد مشکوک با وجود شباهت داشتن با source اشتراک دو متن کم بود؛ یعنی شباهت جاکاردشون کم بود ولی با این وجود به عنوان سرقت ادبی در نظر گرفته می‌شد (با توجه به توضیحات فایل xml آن‌ها) و اگر آستانه را بالا می‌گذاشتیم این نمونه‌های درست حذف می‌شدند و با توجه به این موضوع و تست کردن با ترشهولد بالاتر متوجه شدیم با پایین آوردن ترشهولد نتایج بهتری به دست می‌آید.

مورد دیگری که هست، اگر در این الگوریتم می‌آمدیم شباهت جاکارد جفت اسناد را حساب می‌کردیم می‌شد انتخاب دو از $n$ که پیچیدگی زمانی آن برابر $n^2$ می‌شود که برای ما که در حال کار با دیتاست بزرگی هستیم باعث از دست رفتن زیاد هم زمان هم منابع سخت‌افزار می‌شود. ولی با اجرای این الگوریتم به‌جای بررسی جفت‌جفت فایل‌ها، ما فقط فایل‌هایی که کاندید هستند و باهم شبیه‌اند را بررسی می‌کنیم و این زمان اجرای ما را از $n^2$ خیلی پایین می‌آورد.

\subsection{مسیر دوم: TF-IDF + SimHash}

\begin{itemize}
    \item \textbf{وزن‌دهی TF:} $TF(t,d) = 1 + \log_{10}(\text{count})$ (لگاریتمی)
    \item \textbf{IDF:} $IDF(t) = \log_{10}\left(\dfrac{N}{df(t)}\right)$ (بدون هموارسازی)
    \item \textbf{TF-IDF:} حاصل‌ضرب $TF \times IDF$
    \item \textbf{هش ۶۴ بیتی:} با استفاده از ۸ بایت اول MD5 هر توکن
    \item \textbf{بردار وزن‌دهی‌شده:} برای هر توکن، اگر بیت $i$ هش برابر $1$ باشد، وزن به مؤلفه‌ی $i$ام اضافه و در غیر این‌صورت کم می‌شود.
    \item \textbf{اثرانگشت نهایی:} بیت $i$ برابر $1$ اگر مؤلفه‌ی متناظر $> 0$
    \item \textbf{شباهت:} $1 - \dfrac{\text{فاصله‌ی همینگ}}{\text{تعداد بیت‌ها}}$
\end{itemize}

\section{انتخاب پارامترها (Parameter Choices)}

\begin{center}
\begin{tabular}{|c|c|p{7cm}|}
\hline
\textbf{پارامتر} & \textbf{مقدار} & \textbf{دلیل} \\
\hline
\texttt{shingle\_size (k)} & ۳ & مقدار پیشنهادی در PDF (مشخصات پروژه) قابل تغییر به ۵ نیز هست. \\
\hline
\texttt{num\_perm (MinHash)} & ۱۲۸ & مقدار پیشنهادی در PDF (مشخصات پروژه) قابل تغییر به ۲۵۶ نیز هست. \\
\hline
\texttt{bands, rows} & ۳۲, ۴ & با استفاده از فرمول منحنی S بهینه‌سازی شده است. با این مقادیر، نقطه‌ی گذار (Threshold) حدود ۰٫۴۲ به‌دست می‌آید که برای تشخیص سرقت ادبی بسیار مناسب است. همچنین با توجه به اینکه \texttt{num\_perm = 128} و \texttt{bands × rows = 64 × 2 = 128}، این انتخاب از نظر محاسباتی نیز بهینه است. \\
\hline
\texttt{hash\_prime (LSH)} & $2^{31}-1$ & عدد اول بزرگ برای کاهش برخورد سطل‌ها با حافظه‌ی قابل‌مدیریت \\
\hline
\texttt{hash\_bits (SimHash)} & ۶۴ & مطابق با مشخصات پروژه \\
\hline
\texttt{stemming} & فعال & یکسان‌سازی شکل‌های مختلف کلمات (افزایش بازخوانی) \\
\hline
\texttt{stopword removal} & فعال & حذف کلمات کم‌اطلاعات (افزایش دقت) \\
\hline
\end{tabular}
\end{center}

در پیاده‌سازی نهایی، از \texttt{bands=64} و \texttt{rows=2} استفاده شده است؛ این تغییر برای تنظیم دقیق‌تر نقطه‌ی گذار و بهبود کارایی روی دیتاست صورت گرفته است.

\section{معرفی دیتاست}

از مجموعه‌داده‌ی \textbf{PAN Plagiarism Corpus 2011 (PAN-PC-11)} استفاده شد. این دیتاست شامل دو بخش اصلی است:

\begin{itemize}
    \item \textbf{بخش تشخیص خارجی (External Detection Corpus):} شامل اسناد منبع (\texttt{source-document}) و اسناد مشکوک (\texttt{suspicious-document}) که سرقت‌های دستی (توسط انسان) در آن‌ها جاسازی شده است.
    \item \textbf{فایل‌های XML} در کنار اسناد مشکوک، اطلاعات Ground Truth را به‌صورت تگ \texttt{<feature name="plagiarism">} نگهداری می‌کنند.
\end{itemize}

ویژگی‌های دیتاست مورد استفاده:

\begin{itemize}
    \item تعداد اسناد منبع: حدود ۱۰٬۰۰۰ سند
    \item تعداد اسناد مشکوک: حدود ۱۰٬۰۰۰ سند
    \item نسبت اسناد دارای سرقت ادبی: حدود ۲۰٪
    \item زبان: انگلیسی
\end{itemize}

\section{نتایج و ارزیابی}

پس از اجرای هر دو روش نتایج زیر به‌دست آمد:

\subsection{معیارهای ارزیابی}

\subsubsection*{جدول مقایسه‌ی دقیق معیارها}

\begin{center}
\begin{tabular}{|l|c|c|p{6cm}|}
\hline
\textbf{معیار} & \textbf{MinHash+LSH} & \textbf{SimHash} & \textbf{توضیح} \\
\hline
TP & ۱۳۲ & ۱۱۰ & تشخیص درست سرقت \\
\hline
FP & ۱۴۱ & ۱۱۰٬۸۲۰ & تشخیص نادرست (مثبت کاذب) \\
\hline
FN & ۱۷٬۵۴۲ & ۱۷٬۵۶۴ & از قلم‌افتاده (منفی کاذب) \\
\hline
Precision & ۰٫۴۸۳۵ & ۰٫۰۰۰۹۹ & دقت تشخیص \\
\hline
Recall & ۰٫۰۰۷۵ & ۰٫۰۰۶۲ & بازخوانی (پوشش) \\
\hline
F1-Score & ۰٫۰۱۴۷ & ۰٫۰۰۱۷ & میانگین هارمونیک دقت و بازخوانی \\
\hline
MAP & ۰٫۰۱۱۸ & ۰٫۰۰۴۳ & میانگین دقت میانگین \\
\hline
MRR & ۰٫۰۲۰۰ & ۰٫۰۰۷۹ & میانگین معکوس رتبه‌ی اولین نتیجه‌ی درست \\
\hline
Precision@10 & ۰٫۰۱۹۹ & ۰٫۰۰۲۰ & دقت در ۱۰ نتیجه‌ی برتر \\
\hline
Recall@10 & ۰٫۰۱۱۸ & ۰٫۰۱۰۳ & بازخوانی در ۱۰ نتیجه‌ی برتر \\
\hline
Hit@10 & ۰٫۰۲۰۰ & ۰٫۰۱۹۸ & درصد پرس‌وجوهایی که حداقل یک نتیجه‌ی درست در ۱۰ تای اول دارند \\
\hline
\end{tabular}
\end{center}

\subsubsection*{تحلیل مقایسه‌ای}

\begin{enumerate}
    \item \textbf{دقت (Precision)} \\
    LSH با \textbf{۴۸٫۳۵٪} عملکرد بسیار بهتری نسبت به SimHash با \textbf{۰٫۰۹۹٪} دارد. این اختلاف عمدتاً ناشی از مرحله‌ی \textbf{باندبندی (LSH)} است که تعداد کاندیدها را به شدت کاهش می‌دهد و سپس با محاسبه‌ی \textbf{جاکارد دقیق} روی کاندیدها، تعداد مثبت‌های کاذب را کم می‌کند. در مقابل، SimHash بدون فیلتر کردن، همه‌ی اسناد منبع را با هر سند مشکوک مقایسه می‌کند و در نتیجه ۱۱۰٬۸۲۰ نتیجه‌ی کاذب تولید می‌کند که دقت را به شدت کاهش می‌دهد.

    \item \textbf{بازخوانی (Recall)} \\
    هر دو روش بازخوانی بسیار پایینی دارند (کمتر از ۰٫۸٪). این نشان می‌دهد که اکثر موارد سرقت ادبی در دیتاست (بیش از ۹۹٪) توسط سیستم تشخیص داده نشده‌اند. دلیل اصلی حساسیت شینگل‌های کلمه‌ای به تغییرات واژگانی (مانند تعویض کلمات با مترادف‌ها) و همچنین تنظیم آستانه‌ی نسبتاً بالا (۰٫۲۵) است که باعث می‌شود بسیاری از جفت‌های دارای سرقت، از آستانه عبور نکنند. تفاوت Recall بین دو روش ناچیز است (LSH ۰٫۰۰۷۵ در مقابل SimHash ۰٫۰۰۶۲).

    \item \textbf{نمره‌ی F1} \\
    با وجود دقت بالاتر LSH، به‌دلیل بازخوانی بسیار پایین، F1 هر دو روش کمتر از ۰٫۰۲ است. LSH با F1=۰٫۰۱۴۷ نسبتاً بهتر عمل کرده است.

    \item \textbf{معیارهای رتبه‌بندی (MAP, MRR, Precision@K, Recall@K, Hit@K)} \\
    در همه‌ی این معیارها، LSH برتری محسوسی نسبت به SimHash دارد. برای مثال، MRR در LSH برابر ۰٫۰۲ است در حالی که در SimHash برابر ۰٫۰۰۷۹ است؛ یعنی به‌طور متوسط، اولین نتیجه‌ی درست در LSH در رتبه‌ی ۵۰ (چون $1/0.02=50$) و در SimHash در رتبه‌ی ۱۲۶ قرار دارد. Hit@10 تقریباً در هر دو روش یکسان (حدود ۲٪) است، یعنی از هر ۵۰ سند مشکوک، فقط یک سند حداقل یک منبع درست را در بین ۱۰ نتیجه‌ی برتر دارد.

    \item \textbf{تعداد مثبت‌های کاذب (FP)} \\
    تفاوت فاحش: LSH فقط ۱۴۱ FP دارد، در حالی که SimHash بیش از ۱۱۰ هزار FP تولید کرده است. این نشان می‌دهد که LSH با موفقیت تعداد مقایسه‌ها را از همه‌به‌همه به تعداد کاندیدهای کوچک کاهش داده است.
\end{enumerate}

\subsubsection*{جمع‌بندی ارزیابی}

\begin{itemize}
    \item \textbf{از نظر دقت و کارایی محاسباتی}، روش MinHash+LSH به مراتب برتر از SimHash است و تعداد مثبت‌های کاذب را به طرز چشمگیری کاهش می‌دهد.
    \item \textbf{از نظر بازخوانی}، هر دو روش عملکرد ضعیفی دارند و نیاز به بهبود دارند. پیشنهادهایی برای افزایش بازخوانی:
    \begin{itemize}
        \item کاهش آستانه‌ی شباهت (مثلاً از ۰٫۲۵ به ۰٫۱۵)
        \item استفاده از شینگل‌های کوتاه‌تر (مانند \texttt{k=2}) یا افزودن شینگل‌های کاراکتری
    \end{itemize}
\end{itemize}

\subsection{تحلیل نتایج}

\begin{itemize}
    \item \textbf{دقت (Precision):} MinHash+LSH بالاتر است، زیرا مرحله‌ی LSH کاندیدهای نامرتبط را به‌خوبی فیلتر می‌کند و سپس جاکارد دقیق روی کاندیدها اعمال می‌شود.
    \item \textbf{بازخوانی (Recall):} MinHash+LSH نیز برتری جزئی دارد، چون باندهای متعدد امکان تشخیص حتی اسناد با شباهت متوسط را افزایش می‌دهند.
    \item \textbf{سرعت:} SimHash بسیار سریع‌تر است، زیرا تنها یک اثرانگشت ۶۴-بیتی تولید و با XOR مقایسه می‌شود، در حالی که MinHash+LSH نیاز به تولید امضای ۱۲۸-تایی و جستجوی سطل‌ها دارد.
    \item \textbf{مقایسه‌ی کلی:} اگر سرعت اولویت باشد، SimHash انتخاب مناسبی است. اما برای کاربردهای دقیق‌تر (مانند تشخیص سرقت ادبی در اسناد بلند)، MinHash+LSH دقت بالاتری ارائه می‌دهد.
\end{itemize}

\section{تحلیل خطا (Error Analysis)}

\subsection{مثبت‌های کاذب (False Positives)}

\textbf{مثال:} دو سند با موضوعات کاملاً متفاوت اما با چند واژه‌ی مشترک (مثلاً «computer» و «science») در یک باند LSH قرار می‌گیرند و به‌عنوان کاندید شناخته می‌شوند.

\textbf{دلیل:}
\begin{itemize}
    \item در MinHash+LSH، اشتراک چند شینگل می‌تواند باعث اشتراک در یک باند شود، حتی اگر شباهت کلی پایین باشد.
    \item در SimHash، ممکن است وزن‌های TF-IDF برای کلمات پرتکرار (مانند «data» در متون علمی) باعث نزدیکی اثرانگشت‌ها شود.
\end{itemize}

\textbf{راه‌حل:} افزایش تعداد باندها یا کاهش \texttt{rows} برای افزایش گزینش‌گری، یا تنظیم آستانه‌ی جاکارد در مرحله‌ی تأیید.

\subsection{منفی‌های کاذب (False Negatives)}

\textbf{مثال:} دو سند که به‌شدت بازنویسی شده‌اند (با تعویض گسترده‌ی کلمات) ولی محتوای یکسانی دارند، در هیچ باند مشترکی قرار نمی‌گیرند.

\textbf{دلیل:}
\begin{itemize}
    \item شینگل‌های کلمه‌ای با \texttt{k=3} به تغییرات جزئی حساس هستند. بازنویسی با مترادف‌ها باعث تغییر شینگل‌ها و در نتیجه تغییر امضای MinHash می‌شود.
    \item در SimHash، اگر وزن‌های TF-IDF به‌خاطر کورپوس کوچک نتوانند تمایز ایجاد کنند، اثرانگشت‌ها مشابه می‌شوند.
\end{itemize}

\textbf{راه‌حل:} استفاده از شینگل‌های کوتاه‌تر (\texttt{k=2}) یا اضافه کردن مرحله‌ی تشخیص مترادف (که در این پروژه پیاده‌سازی نشده است).

\subsection{محدودیت‌های کلی}

\begin{itemize}
    \item \textbf{وابستگی به کورپوس:} IDF در SimHash به دقت بر روی دیتاست مرجع محاسبه می‌شود. در کورپوس‌های کوچک، تمایز کلمات کاهش می‌یابد (مشکل در \texttt{compare} دو فایل).
    \item \textbf{عدم پشتیبانی از متون فارسی:} تنها متون انگلیسی پردازش می‌شوند.
    \item \textbf{هش ۶۴ بیتی در SimHash:} احتمال برخورد هش برای کلمات بسیار زیاد در دیتاست‌های غول‌پیکر وجود دارد (هرچند در عمل نادر است).
\end{itemize}

\section{نتیجه‌گیری (Conclusion)}

پروژه‌ی حاضر یک سامانه‌ی کامل و ماژولار برای تشخیص سرقت ادبی و اسناد مشابه ارائه می‌دهد که با دو رویکرد متفاوت (MinHash+LSH و SimHash) پیاده‌سازی شده است.

\begin{itemize}
    \item \textbf{MinHash+LSH} دقت و بازخوانی بالاتری دارد ولی کندتر است.
    \item \textbf{SimHash} بسیار سریع‌تر است و برای جستجوی در مقیاس بزرگ مناسب‌تر است، اما با افت جزئی دقت همراه است.
\end{itemize}

انتخاب روش نهایی به نیاز برنامه‌ی کاربردی (سرعت در مقابل دقت) بستگی دارد. همچنین پیش‌پردازش دقیق و انتخاب صحیح پارامترها نقش کلیدی در عملکرد نهایی ایفا می‌کنند.

\section{پیوند به کد منبع و مستندات}

\begin{itemize}
    \item \textbf{مخزن GitHub:} \href{https://github.com/your-username/semantic-plagiarism-engine}{semantic-plagiarism-engine}
    \item \textbf{گزارش کامل (LaTeX):} \texttt{docs/project\_spec.tex}
    \item \textbf{نتایج و فایل‌های خروجی:} \texttt{outputs/metrics.csv}, \texttt{outputs/candidates.csv}
\end{itemize}

\end{document}